\documentclass[11pt,a4paper]{article}
\usepackage[utf8]{inputenc}
\usepackage[T1]{fontenc}
\usepackage[margin=25mm]{geometry}
\usepackage{xcolor}
\usepackage{graphicx}
\usepackage{hyperref}
\definecolor{forest}{HTML}{183D3B}
\definecolor{copper}{HTML}{8A582F}
\hypersetup{colorlinks=true,linkcolor=forest,urlcolor=forest,pdftitle={TerraSync: documento técnico},pdfauthor={404 Gravity Not Found}}
\setlength{\parindent}{0pt}
\setlength{\parskip}{7pt}
\renewcommand{\abstractname}{Resumen}
\renewcommand{\contentsname}{Contenido}
\newcommand{\proyecto}{TerraSync}
\title{\textcolor{forest}{\Huge\proyecto}\\[8pt]Atlas de observación terrestre de México\\[8pt]\large Documento técnico}
\author{Equipo 404 Gravity Not Found}
\date{15 de noviembre de 2026}
\begin{document}
\maketitle
\begin{abstract}
\proyecto{} integra un catálogo geoespacial mensual con un explorador de México, una tabla de observaciones y recursos para interpretar la evidencia. Esta entrega conecta React y MapLibre con Django REST Framework, PostgreSQL y archivos Parquet. Añade selección temporal mediante slider, filtros coherentes por entidad y fenómeno, exportación Excel y una interfaz en español e inglés. El sistema se ejecuta localmente; se incluye un perfil Docker con Nginx y Gunicorn para ensayar un despliegue sin publicación externa. La validación descrita comprueba comportamiento informático y no constituye validación científica independiente del catálogo.
\end{abstract}
\tableofcontents
\newpage
\section{Problema y alcance}
Los registros de observación terrestre requieren contexto espacial, temporal y metodológico. Un punto en un mapa puede representar una observación, un indicador indirecto o una salida de modelo; presentarlos sin distinguir su evidencia puede inducir interpretaciones incorrectas.

El objetivo de \proyecto{} es permitir que una persona seleccione un estado, recorra el registro mensual, examine las fuentes y descargue exactamente el subconjunto consultado. La interfaz tiene cuatro rutas: Inicio, Explorar, Observaciones y Aprende. La portada utiliza una representación ilustrativa de la Tierra y de una órbita; esa animación no representa telemetría de NISAR.

El catálogo incluido contiene 3,290 registros espaciales, 203 meses y cobertura de las 32 entidades. Estas cifras describen el artefacto recibido, no una cobertura uniforme de todos los fenómenos. La interfaz no consulta una transmisión satelital en vivo y no proporciona alertas oficiales ni imágenes SAR procesadas para descarga.

\section{Arquitectura}
\begin{center}
\fbox{React + TypeScript}\quad$\longrightarrow$\quad\fbox{Django REST API}\\[10pt]
\hspace*{55mm}$\downarrow$\\[6pt]
\hspace*{35mm}\fbox{PostgreSQL + catálogo Parquet}
\end{center}
\subsection{Frontend}
React organiza las páginas y controles. TypeScript verifica los contratos estáticos y el cliente valida la estructura del GeoJSON recibido antes de dibujar puntos. Vite compila el proyecto; en desarrollo reenvía el prefijo \texttt{/api/} a Django. MapLibre mantiene fuentes separadas para estados, municipios, etiquetas e indicadores. Las colecciones municipales se cargan únicamente al seleccionar una entidad.

El explorador usa filas de CSS Grid independientes para cabecera, filtros, lienzo, estado de consulta y línea del tiempo. La navegación ocupa una fila adicional. Esta separación evita que el historial o los controles oculten contenido del mapa. En pantallas pequeñas el panel territorial sustituye temporalmente el lienzo mientras está abierto; al cerrarlo, el mapa recupera su espacio mediante su observador de tamaño.

El slider nativo representa índices enteros sobre los periodos devueltos por la API. Así, sus pasos corresponden a meses realmente disponibles para los filtros, aunque existan huecos en la serie. Permite arrastre, teclado, avance por mes, selector de periodo y reproducción con pausa. El histograma asociado representa la severidad máxima de cada periodo; no representa probabilidades calibradas.

\subsection{Backend y almacenamiento}
Django REST Framework sirve colecciones GeoJSON, resúmenes temporales y archivos Excel. PostgreSQL almacena los eventos espaciales; el servicio de catálogo sincroniza los registros del Parquet incluido. El volumen del catálogo se monta en modo de solo lectura. La caché del proceso reduce lecturas posteriores; reemplazar archivos de catálogo requiere reiniciar los procesos para invalidar sus cachés.

Se conserva el directorio del backend anterior como referencia. El perfil activo utiliza la versión de \texttt{404-Gravity-Not-Found-integration}, que implementa los servicios científicos y el historial. No deben iniciarse ambos backends sobre el mismo puerto.

\section{Contrato de datos y consultas}
Cada evento contiene geometría puntual y propiedades: identificador, fecha, periodo, tipo de fenómeno, entidad, nombre, severidad, radio, evidencia, fuente y, cuando existen, puntuación y probabilidad. Los nombres y etiquetas científicas originales se conservan aunque el idioma de la interfaz cambie.

\begin{tabular}{ll}
\textbf{Ruta GET} & \textbf{Resultado}\\ \hline
\texttt{/api/predictions/} & Colección GeoJSON del periodo\\
\texttt{/api/timemap/} & Periodos, conteos y máximos\\
\texttt{/api/catalog/status/} & Estado y cobertura del catálogo\\
\texttt{/api/export/} & Libro Excel del subconjunto seleccionado\\
\end{tabular}

Los parámetros principales son \texttt{period=YYYY-MM}, \texttt{state=01..32} y \texttt{type}. Los tipos de amenaza son \texttt{sismo}, \texttt{huracan}, \texttt{sequia}, \texttt{inundacion} e \texttt{incendio}. Un ejemplo reproducible es:
\begin{quote}\small\texttt{/api/predictions/?period=2026-11\&state=12\&type=incendio}\end{quote}
El estado seleccionado por clic y el selector de entidades comparten la misma variable. Esa variable se envía al mapa, al historial, al enlace de Observaciones y a la exportación. Si una consulta filtrada no tiene resultados, se devuelve una colección vacía: no se sustituye por registros agregados nacionales.

Las solicitudes del frontend usan \texttt{AbortController}. Cambiar de mes o entidad cancela la solicitud anterior; una respuesta obsoleta no debe cambiar el estado de la consulta vigente. Los controles informan de carga, error y ausencia de datos, y permiten reintentar.

\section{Evidencia y límites de interpretación}
Las etiquetas se preservan desde el catálogo de origen:
\begin{description}
\item[observed] Registro marcado como observado por el artefacto de entrada.
\item[supervised] Salida etiquetada como supervisada en el catálogo.
\item[proxy] Indicador indirecto que requiere contexto y corroboración.
\item[no\_confirmable] Evidencia insuficiente para confirmar el fenómeno.
\end{description}
Una etiqueta no constituye una auditoría independiente de su procedencia. Una severidad numérica no debe interpretarse automáticamente como probabilidad de ocurrencia. La ausencia de registros tampoco demuestra ausencia de fenómenos. Antes de cualquier uso operativo deben revisarse localización, cobertura temporal, construcción de variables, calibración y fuentes originales. No se reportan nuevas métricas de precisión científica en esta entrega.

\section{Interfaz, idioma y accesibilidad}
La dirección visual utiliza papel cálido, tinta verde y acentos cobre. IBM Plex Sans se incluye localmente con su licencia SIL Open Font License; los títulos editoriales usan Georgia. Los colores de fenómenos se conservan como códigos semánticos.

El selector ES/EN está en la esquina superior derecha y mantiene la preferencia en almacenamiento local cuando el navegador lo permite. Actualiza el atributo \texttt{lang} del documento. Nombres propios, nombres de estados, fuentes y valores científicos no se traducen automáticamente.

El tutorial es un diálogo nativo modal de cuatro pasos con foco y cierre explícito. Aprende incluye lecciones desplegables y una actividad con retroalimentación. Observaciones muestra una tabla con filtros, estados vacíos, errores y exportación. Se retiran textos de preparación obsoletos y pies redundantes; se conservan la atribución cartográfica y las advertencias de interpretación necesarias.

\section{Exportación Excel}
El endpoint de exportación reutiliza la misma selección del endpoint GeoJSON. Genera un archivo \texttt{.xlsx} mediante el contenedor ZIP y XML de Office Open XML, sin dependencias de hoja de cálculo adicionales. Incluye una hoja con encabezados, fila congelada y autofiltro.

Se exportan fecha, periodo, entidad, fenómeno, nombre, severidad, evidencia, observación, puntuación, probabilidad, fuente y coordenadas. Los valores textuales se escriben como cadenas en línea, incluso si comienzan por el signo igual; no se convierten en fórmulas ejecutables. Las columnas numéricas se preservan como números. Una selección vacía conserva los encabezados y no introduce filas ficticias.

\section{Ejecución y despliegue local}
\subsection{Desarrollo}
Desde la raíz del repositorio:
\begin{verbatim}
pnpm stack:up
pnpm dev --host 127.0.0.1 --port 5173 --strictPort
\end{verbatim}
El frontend se abre en \url{http://127.0.0.1:5173/explorar}; la API se expone en \url{http://127.0.0.1:8000}. PostgreSQL permanece dentro de la red Docker.

\subsection{Perfil de despliegue}
\texttt{compose.deploy.yaml} construye React en una etapa Node y sirve el resultado con Nginx. Django se ejecuta mediante Gunicorn con depuración desactivada. El tráfico \texttt{/api/} se reenvía internamente al backend, por lo que no necesita un origen público distinto. PostgreSQL tiene un volumen separado del entorno de desarrollo.
\begin{verbatim}
docker compose --env-file .env.deploy \
  -f compose.deploy.yaml up --build -d --wait
\end{verbatim}
La escucha predeterminada es únicamente \url{http://127.0.0.1:8080}. Las variables \texttt{SECRET\_KEY} y \texttt{DATABASE\_PASSWORD} son obligatorias y se mantienen en archivos ignorados por Git. Este trabajo no publica el sistema en Internet. Un despliegue futuro requiere dominio, TLS, revisión de orígenes, copias de seguridad, observabilidad y política operativa.

\section{Validación y limitaciones conocidas}
Se ejecutan la comprobación TypeScript, el linter, la compilación y las pruebas del catálogo en una base temporal. Se añadieron regresiones para impedir el retorno nacional ante un filtro estatal vacío y para verificar que los textos exportados no se escriben como fórmulas.

Las comprobaciones de navegador cubren 1440, 1024, 768, 390 y 320 píxeles, el cambio de mes por slider, el filtro territorial, el selector de idioma, el tutorial, la tabla y la actividad educativa. Las filas del mapa, historial y navegación se comprueban por sus rectángulos, además de revisar capturas. No se certifican todas las combinaciones de navegador, rendimiento de GPU física ni accesibilidad completa con lectores de pantalla.

El build conserva avisos de tamaño de los motores cartográficos y 3D. El fondo topográfico requiere Internet; la cartografía local mantiene su disponibilidad aunque falle ese proveedor. Las cachés y límites por proceso requieren una solución compartida antes de escalar horizontalmente. Los nombres de integrantes siguen pendientes de confirmación; se configuran en \texttt{src/configuracion/proyecto.ts}. El isotipo incluido es una propuesta vectorial editable.

\section{Fuentes y repositorio}
\begin{itemize}
\item Código: \url{https://github.com/GalekTato/404-Gravity-Not-Found}
\item NASA: \href{https://science.nasa.gov/mission/nisar/get-to-know-sar/}{Get to know SAR: fundamentos del radar de apertura sintética}.
\item INEGI: \href{https://www.inegi.org.mx/servicios/catalogounico/}{Catálogo Único de Claves de Áreas Geoestadísticas}.
\item Procedencia de cartografía: \path{public/geojson/README.md} y \path{procedencia.json}.
\item Catálogo científico: archivos Parquet y manifiestos.\\\path{404-Gravity-Not-Found-integration/data/}
\item Fuentes tipográficas: licencia incluida en \path{public/fonts/OFL.txt}.
\end{itemize}
\end{document}
